Metal-poor stars preserve a fossil record of the Milky Way’s early assembly, but extinction and crowding have left the inner Galactic midplane largely inaccessible to optical searches. We propose a 60 square degree pathfinder imaging survey with Blanco/NEWFIRM using the H1 and H2 intermediate-band filters. Synthetic photometry predicts that H1-H2, combined with existing broad-band J-Ks photometry, can distinguish upper red-giant-branch stars with [Fe/H] < -1 from the much more numerous metal-rich red-clump and disk populations. In two nights, we will obtain 60 s per filter over 306 pointings, reaching S/N ≈ 140 per band and σ(H1-H2) ≈ 0.011 mag at the Hidden Galaxy Explorer (HGE) targeting limit of H = 15.5. HGE will begin deep near-infrared spectroscopy of the dust-obscured far Milky Way in mid-2027. Existing APOGEE spectroscopy will anchor the bright color-metallicity relation, while early HGE spectra will validate the selection at fainter magnitudes. The resulting catalog will provide a ranked list of metal-poor candidates for early HGE operations and measure how the selection performs across representative variations in extinction and crowding. This program will establish the strategy for extending NEWFIRM imaging across the broader HGE footprint and deliver a uniform near-infrared legacy data set for the inner Galactic midplane.